\centerline{\bf \Huge Геометрия 1}
\centerline{\bf \Huge Точка Микеля}

\begin{flushright}
        \scriptsize{Все решения должны быть в направленных углах!}
\end{flushright}

\problem{1}
Четыре прямые $a, b, c$ и $d$ общего положения высекают на плоскости выпуклый четырехугольник $A B C D$.

(a) Докажите, что противоположные стороны этого четырехугольника видны из точки Микеля $M$ этой четверки прямых под равными углами.

(b) Докажите, что $M A \cdot M C=M B \cdot M D$.

\problem{2}
Четыре прямые общего положения высекают на плоскости четыре треугольника. Докажите, что центры описанных около этих треугольников окружностей лежат на одной окружности, проходящей через точку Микеля этих прямых.

\problem{3}
У четырехугольника $A B C D$ нет параллельных сторон. Прямая $\ell$ пересекает прямые, содержащие стороны четырехугольника, в четырех точках. Восстановим в каждой из этих точек перпендикуляр к соответствующей стороне четырехугольника. Докажите, что прямую $\ell$ можно выбрать так, чтобы эти перпендикуляры пересеклись в одной точке.

\problem{4}
5-окружность Микеля. Даны пять прямых общего положения. Рассмотрим все пять точек Микеля всевозможных четверок этих прямых. Докажите, что получившееся пять точек будут лежать на одной окружности или прямой.

\problem{5}
Дан четырехугольник $A B C D$, вписанный в окружность $\omega$ с центром в точке $O$. Прямые $A B$ и $C D$ пересекаются в точке $P$, а прямые $B C$ и $A D$ - в точке $Q$, а прямые $A C$ и $B D$ - в точке $R$.
\begin{itemize}
    \item[(a)] Докажите, что точка Микеля $M$ четверки прямых, содержащих стороны четырехугольника, лежит на прямой $P Q$.

    \item[(b)] Докажите, что четырехугольники $M A O C$ и $B O D M$ - вписанные.

    \item[(c)] Докажите, что $M O \perp P Q$.

    \item[(d)] Покажите, что точка $R$ лежит на прямой $O M$ и выведите из этого, что $O R \perp P Q$.

    \item[(e)] Теорема. Точка $O$ является ортоцентром треугольника $P Q R$

    \item[(f)] Обозначим за $X$ и $Y$ основания касательных, проведенных из точки $P$ к $\omega$. Докажите, что прямая $X Y$ проходит через точку $Q$.
\end{itemize}


\problem{6}
На сторонах $A B$ и $A C$ треугольника $A B C$ выбираются такие точки $P$ и $Q$, что четырехугольник $B P Q C$ вписан в окружность с центром в точке $O$. Диагонали четырехугольника пересекаются в точке $S$. Докажите, что прямая $O S$ проходит через одну из общих точек окружностей ( $B A C$ ) и $(P A Q)$.